%!TeX encoding = UTF-8 Unicode
%!TeX spellcheck = en-UK
% !TeX root = Perrinet26mnesis.tex
% -*- coding: UTF-8; -*-
%%%-----------------------------------------------------------------
\begin{figure}[!t] 
  \centering
  % the panels are printed at fig_width/3 by notebook 20, hence scaled by 3/4 here: use fig_width/4 at its next run
  \includegraphics[width=0.25\linewidth]{scan_num_delay.pdf}\hfill
  \includegraphics[width=0.25\linewidth]{scan_N_neuron.pdf}\hfill
  \includegraphics[width=0.25\linewidth]{scan_N_time.pdf}\hfill
  \includegraphics[width=0.25\linewidth]{scan_N_pattern.pdf}
  % TODO \includegraphics[width=0.9\linewidth]{scan_p_A.pdf}
  \caption{%
    Effect of some key parameters on working-memory performance ($F_1$ score over $10$ independent runs, each training a new network on new patterns drawn from its own random seed, a given seed sharing its patterns across all values of a parameter; bars show the mean over these runs, and whiskers span the $5\%$ to $95\%$ quantiles of a beta distribution fitted to them). %
    \emph{From left to right:}
    The $F_1$ score rises steeply with the maximum delay $D$ and with the number of neurons $N$: both increase the mean number $p_A N D$ of spikes in a context window, and thus decrease the probability that a motif contains a gap of $D$ silent steps, at which its replay stops. %
    It decreases slowly with the pattern duration $T$, as a longer motif is more likely to contain such a gap. %
    It does not depend on the number of stored motifs $M$ over the scanned range, up to $M = 26$.
  }\label{fig:scan}
\end{figure}
%%%-----------------------------------------------------------------
